\chapter{Testing and Evaluation}
\label{chap:10}

\section{Test Strategy}
\label{sec:10.1}

Three levels, chosen so that each one tests something the others cannot reach and so that the cheapest level carries the most cases.

{\small
\setlength{\LTpre}{6pt}\setlength{\LTpost}{8pt}
\begin{longtable}{>{\raggedright\arraybackslash}p{\dimexpr 0.213\textwidth-2\tabcolsep\relax}>{\raggedright\arraybackslash}p{\dimexpr 0.345\textwidth-2\tabcolsep\relax}>{\raggedright\arraybackslash}p{\dimexpr 0.442\textwidth-2\tabcolsep\relax}}
\caption{The three levels of testing, and what each one answers}\\
\toprule
Level & Runs on & Answers \\
\midrule
\endfirsthead
\toprule
Level & Runs on & Answers \\
\midrule
\endhead
\bottomrule
\endfoot
\bottomrule
\endlastfoot
Unit and integration & A desktop JVM, no device & Does the evaluator do what the specification says? \\
Replay & A desktop JVM, recorded video & Does the specification describe the movement? \\
Device & The target handset & Does the whole path work, at the rate it must? \\
\end{longtable}
}

The split is possible because of one architectural property: the assessment module has no Android dependency, so the layer carrying almost all of the logic can be tested without an emulator, a device or a camera. 5.4 records how that boundary is enforced.

The two questions in the middle column are genuinely different, and conflating them is the easiest mistake available here. A suite that passes establishes that a threshold is applied correctly. Whether the threshold is the right number for the movement is a question only recorded video can answer, which is what the replay level exists for.

\textbf{Determinism is what makes the bottom level possible.} NFR-A4 requires identical landmark input to produce identical output. Every unit test depends on it: each replays a fixed sequence and asserts an exact verdict. If the evaluator carried any hidden state or any dependence on wall-clock time, the suite could not be written at all.

\section{Unit and Integration Tests}
\label{sec:10.2}

\textbf{85 tests in the assessment module and 9 in the client, 0 failures.}

{\small
\setlength{\LTpre}{6pt}\setlength{\LTpost}{8pt}
\begin{longtable}{>{\raggedright\arraybackslash}p{\dimexpr 0.318\textwidth-2\tabcolsep\relax}>{\raggedright\arraybackslash}p{\dimexpr 0.682\textwidth-2\tabcolsep\relax}}
\caption{What the unit and integration tests assert}\\
\toprule
Concern & What is asserted \\
\midrule
\endfirsthead
\toprule
Concern & What is asserted \\
\midrule
\endhead
\bottomrule
\endfoot
\bottomrule
\endlastfoot
Specification loading & All eleven validation rules reject what they are meant to reject, and the failure is named rather than generic \\
Measurement vocabulary & Each function returns the expected value on known geometry, including when a landmark is missing \\
Segmentation & Hysteresis, phase resolution, the stillness band, the duration bounds, and the hold machine \\
Session behaviour & Counting, per-repetition and continuous evaluation, the per-repetition buffer emptying at each boundary \\
Trigger policy & Confidence floor, minimum gap, per-set cap, consecutive-cue limit, and priority between simultaneous faults \\
Tempo independence & The same movement performed fast and slow produces the same verdicts \\
\end{longtable}
}

Traced to the requirements they serve:

{\small
\setlength{\LTpre}{6pt}\setlength{\LTpost}{8pt}
\begin{longtable}{>{\raggedright\arraybackslash}p{\dimexpr 0.500\textwidth-2\tabcolsep\relax}>{\raggedright\arraybackslash}p{\dimexpr 0.500\textwidth-2\tabcolsep\relax}}
\caption{Requirements traced to the tests that cover them}\\
\toprule
Requirement & Tested by \\
\midrule
\endfirsthead
\toprule
Requirement & Tested by \\
\midrule
\endhead
\bottomrule
\endfoot
\bottomrule
\endlastfoot
FR-05 repetition segmentation and counting & Segmentation and session tests \\
FR-06 fault detection per the active specification & Session tests, over all three specifications \\
FR-08 the trigger policy governs speech and silence & Trigger policy tests \\
FR-19 exercises load from files with no code change & The suite runs all three specifications through one evaluator \\
FR-23 an invalid specification is rejected at load & Specification loading tests, one per rule \\
FR-24 every declared fault resolves to an audio asset & Specification loading tests \\
FR-25 every fault carries a provenance record & Specification loading tests \\
NFR-A4 identical input, identical output & Every test in the module, by construction \\
NFR-E3 validation failures are named & Specification loading tests \\
\end{longtable}
}

\textbf{The tests load the real specification files rather than copies.} A specification edited into a state the schema forbids fails this build. This is the only arrangement under which the rules the format claims to enforce stay true of the files that actually ship. It is why FR-25 is a property of the format and not a convention.

\textbf{Tempo independence has a test of its own} because it is the failure mode the segmenter is most prone to. A rule evaluated over a phase will see more frames when the trainee moves slowly, and a confirmation window counted in frames therefore behaves differently at different speeds. The test performs the same movement at two tempos and asserts the same verdicts.

\section{Device Testing}
\label{sec:10.3}

\textbf{Fourteen sets were performed on the target handset with the application running}, across the development of the client. They establish the parts of the system that no desktop test can reach: that the camera path sustains its rate, that the pose engine runs on the GPU delegate under real conditions, and that a confirmed fault reaches the speaker.

{\small
\setlength{\LTpre}{6pt}\setlength{\LTpost}{8pt}
\begin{longtable}{>{\raggedright\arraybackslash}p{\dimexpr 0.500\textwidth-2\tabcolsep\relax}>{\raggedright\arraybackslash}p{\dimexpr 0.500\textwidth-2\tabcolsep\relax}}
\caption{What device testing established, and what it could not}\\
\toprule
Property & Result \\
\midrule
\endfirsthead
\toprule
Property & Result \\
\midrule
\endhead
\bottomrule
\endfoot
\bottomrule
\endlastfoot
Sustained frame rate during a set & 23.91 to 29.99 fps across the fourteen sets \\
Pose rate within a set & up to 0.993 of frames carrying a pose \\
Cues reaching audio & 18 cues spoken across the sets, with the path timed on each \\
Confirmation to audio dispatch, current build & 12 to 20 ms, median 16 \\
\end{longtable}
}

\textbf{Every set in that range clears the NFR-P2 floor of 22 fps}, including the slowest.

Device testing also found the two defects that only hardware could expose, and both are recorded in the chapters that own them. A trigger-policy control that silenced cues entirely while every other indicator looked healthy, which 6.5 describes. And a confidence floor carried unchanged across an engine change, where the two engines report scores on different scales, which 6.5.1 describes. Neither is visible to a unit test, because in both cases the evaluator was behaving exactly as specified.

\section{Evaluation Method}
\label{sec:10.4}

Fault detection accuracy is evaluated against recorded video, scored by the replay harness in 9.3, with the ground truth established two ways.

\textbf{By construction, for the scripted sets.} Each of the ten sets in 9.3.1 was performed to one instruction, so the correct answer is known before the system runs. Repetition counts were confirmed independently of the segmenter, by reading excursions in the driver trace.

\textbf{By an independent expert, for the evaluation pass.} A certified trainer has reviewed the recordings and confirmed that each performs the fault its set is named for. The labelling pass below is a stronger thing and has not been run; the protocol is blind in a specific sense:

\begin{itemize}
\item The trainer sees the video and nothing else. Not the system's output, not the author's labels, and not the fault set the system implements.
\item They are asked what they observe, in their own words, rather than asked to confirm or reject a list.
\item Their labels are compared against the system's afterwards, and agreement is reported as a figure rather than asserted.
\end{itemize}

\textbf{The order matters and cannot be repeated.} Once a labeller has seen the system's output they cannot un-see it, and every judgement afterwards is anchored. This is the one step in the whole evaluation that has exactly one attempt, which is why the protocol fixes it in writing before any footage is sent.

The metrics are those in 3.4.3: precision per fault against NFR-A1's 0.90, recall per fault against NFR-A2's 0.70, and repetition count error against NFR-A3's allowance of one. 9.4 reports all three against the scripted ground truth. The trainer pass replaces the scripted labels with expert ones and reports the same metrics again; what the trainers set is what is reported.

\textbf{Precision is weighted above recall deliberately}, and the gap between the two targets is the requirement rather than an accident of two separately chosen numbers. A missed fault costs one uncorrected repetition. A false cue costs the trainee's confidence in the product, and 3.4.3 records that as the most damaging failure this system has.

\section{Limits of the Current Testing}
\label{sec:10.5}

\textbf{A green suite proves consistency with the specification, not correctness of it.} Every unit test asserts that the evaluator applies a declared threshold as declared. None of them can say whether that threshold is the right number for a human movement, and no arrangement of unit tests could. That question belongs entirely to the recorded sets and to the trainer pass.

\textbf{The test fixtures are synthetic.} They are generated landmark sequences with known geometry, so they exercise the logic cleanly and they do not contain the thing real data contains: noise correlated with the movement itself. A landmark is noisiest exactly where the body occludes it, which is at the bottom of a squat, in the phase where the depth rule is evaluated. The fixtures cannot reproduce that, and the recorded sets are the only place it appears.

\textbf{The evaluation set is one subject in one session.} 9.5 states this in full. The consequence for this chapter is that no figure in 9.4 supports a claim about a second person, and nothing in the testing method would reveal a threshold that happens to suit this performer's proportions.

\textbf{One requirement has no test that closes it.} NFR-P1 is bounded below rather than measured, because the platform path from the play call to the first audible sample is not instrumented. It is stated here rather than left to be discovered. NFR-P4 was closed by the ten-minute run in 9.4.1, which also gives NFR-A5 its evidence; that run carries a caveat of its own about the subject, stated where the figures are.

\textbf{The server has no tests because it has no implementation.} FR-16 to FR-18 describe synchronisation behaviour whose contract is fixed and whose code does not exist in this phase. 11.3 reports that against the objective it serves rather than counting it as passing.

\textbf{Device testing is a sample rather than a regime.} Fourteen sets across the development of the client is enough to establish that the path works and to find two defects that unit tests could not. It is not a stability programme, and no claim is made about behaviour over repeated sessions, across reboots, or on a device other than the one in 4.1.
